\scanpage{66}
``correspondant'' est une isogénie, ou encore, a un noyau fini. Pour bien faire, \deleted{il faudrait}\inserted{on peut} prouver mieux, savoir qu'il correspond à une \uline{polarisation} de $B_X$ (compte tenu de $A_X\simeq B_X^*$).

\deleted{Or cet hom. peut se décrire de proche en proche comme composé d'hom \deleted{de restriction},}

\diagram{\begin{tikzcd}[ampersand replacement=\&,column sep=tiny,row sep=small,execute at end picture={\draw[thin] (current bounding box.south west) -- (current bounding box.north east);}]
\deleted{B_X=B_{X_0}}\arrow[r,"{\deleted{\rho_0}}"] \& \deleted{B_{X_1}}\arrow[r,"{\deleted{\rho_1}}"] \& \deleted{B_{X_2}} \& \deleted{\ldots}\arrow[r,"{\deleted{\rho_{n-3}}}"] \& \deleted{B_{X_{n-2}}}\arrow[r,"{\deleted{\sigma}}"] \& \deleted{A_{X_{\unsure{n-2}}}}\arrow[r,"{\deleted{i}}"] \& \deleted{\unsure{\vert|1}}\\
\& \& \& \& {} \arrow[r] \& \deleted{B_C\simeq A_C}\arrow[r] \& {}
\end{tikzcd}}
\diagram{\begin{tikzcd}[ampersand replacement=\&,execute at end picture={\draw[thin] (current bounding box.north west) -- (current bounding box.south east);}]
\deleted{B_X}\arrow[r,"{\deleted{\rho}}"] \& \deleted{B_Y}\arrow[r,"{\deleted{\lambda}}"] \&\deleted{A_Y}\arrow[r,"{\deleted{\sigma}}"] \&\deleted{A_X}
\end{tikzcd}}

\deleted{où $Y$ est défini \deleted{par} comme l'intersection de $n-2$ \deleted{\unsure{hypersup|hyperplans}} diviseurs ``généraux'' correspondant à $L_1,\ldots,L_{n-2}$ (supposés \uline{très} amples, ce qui est loisible) $\rho$ est la restriction, $\sigma$ l'injection, et $\lambda$ est défini comme $u$, en termes de $L_{n-1}\mid Y$ sur $Y$. On sait \inserted{(de proche en proche\ldots)} que $\rho,\sigma$ sont des isogénies, \deleted{et} (d'ailleurs transposées \deleted{\unsure{sg}} l'une de l'autre),}

Or une fonctorialité évidente en $X,Y$ permet d'interpréter de la façon suivante l'hom $B_X\unsure{\mathord{-}|\to}A_Y$ défini dans 2.1., lorsque $Z$ est défini par un morphisme
% archive sheet number: 65

